\section{Getting started}\label{sec:getting-started}

\subsection{Building \code{qtrocket-cli}}

QtRocket builds with CMake presets; Qt6 is the only system dependency (and the
CLI itself does not use Qt --- Qt is needed to build the sibling GUI targets).
Everything else is fetched and pinned by CMake, so the first configure is slow
once:

\begin{minted}{text}
cmake --preset release
cmake --build --preset release
\end{minted}

The binary lands at \code{build/cli/qtrocket-cli}. The examples in this guide
assume it is on your \code{PATH} or invoked by that path.

\subsection{Four ways to run it}\label{sec:invocation}

\begin{enumerate}
\item \textbf{Interactively} --- run \code{qtrocket-cli} with no arguments and
      type commands at it. It reads lines from standard input until
      end-of-file or a \repl{quit}.
\item \textbf{From a pipe} --- \code{printf 'help\textbackslash{}n' | qtrocket-cli}.
      Same reader, no terminal required.
\item \textbf{From a script file} --- \code{qtrocket-cli flight.cli} runs the
      commands in the file. This is how the tutorial scripts are run.
\item \textbf{One-liner} --- \code{qtrocket-cli -c "listmotors"} runs a single
      command string in a fresh process. \code{--help} and \code{--version}
      do what you expect.
\end{enumerate}

\begin{warning}[State does not survive between processes]
Each invocation is a fresh session. \code{qtrocket-cli -c "setmotor B6"}
followed by \code{qtrocket-cli -c "launch"} does \emph{not} fly a B6 --- the
second process never saw the first command. Put multi-step work in one script
file, one pipe, or one interactive session; use \code{.qrd} design files
(\cref{sec:tut5}) to carry work \emph{between} sessions.
\end{warning}

\subsection{The output protocol and exit codes}\label{sec:protocol}

The CLI is built to be read by both humans and scripts. Every command answers
with a line beginning \code{OK <cmd>} on success or \code{ERR} on failure;
data payloads (motor listings, part trees, state samples) usually follow the
\code{OK} header unprefixed. Two exceptions keep parsers honest: \repl{help}
prints its \code{\#}-prefixed listing \emph{before} its terminating
\code{OK help}, and comment or blank lines produce no output at all. Two more prefixes appear only around flights:
\code{CSV <path>} names the trajectory file \repl{launch} wrote, and
\code{WARN: run aborted (...)} flags a flight that ended for a safety reason
rather than a normal landing. Comment lines start with \code{\#}; the log
system is pinned to errors-only with an \code{[ERROR]} prefix so log noise can
never masquerade as command output.

The process exit code summarizes the whole session: \textbf{0} if every command
succeeded, \textbf{1} if any command reported \code{ERR}, and \textbf{2} for a
bad invocation (unreadable script file, malformed \code{-c}). A failing
command never kills the session --- the REPL reports it and keeps reading ---
but the exit code remembers. \Cref{sec:tut5} turns this into one-line
regression guards.

\subsection{A ninety-second first session}

\begin{console}
\muted{\# QtRocket CLI - type 'help' for commands, 'quit' to exit.}
\uin{help}
\muted{\# commands:}
\muted{\#   loadmotors <file.rse|file.eng>   import a motor database}
\muted{\#   savedb <file.qmd>       save the motor database to a file}
\textcolor{qrGray}{\textit{[40 more help lines --- the full list is \cref{sec:reference}]}}
\ok help
\uin{launch}
\err launch: no motor set (use loadmotors + setmotor first)
\uin{quit}
\ok bye
\end{console}

That \err{} line is the CLI teaching you its own rules: a flight needs a motor,
and motors come from a database you load. That is exactly where
Tutorial~1 begins.

\begin{tip}[If you get lost]
\repl{help} prints every command with a one-line summary, and \repl{status}
(\cref{sec:reference}) shows the current session configuration --- motor,
masses, atmosphere, integrator, timestep --- at any moment.
\end{tip}
